\section{Prior Accomplishments}

As a Ph.D. student at Stanford University, I focused my research
on managing continuous media data.
During my final three years at Stanford, I 
worked on three different
research projects: media servers, media receivers and image databases.
These three areas 
knit nicely together to provide a strong foundation 
for my advanced research in interactive TV.

\subsection{Media Servers}

To deliver a large volume of continuos media data
(i.e., video and audio) from a media
server to a large number of simultaneous clients poses at least
three challenges. First, the system resources
at the server side must be allocated 
carefully to maximize throughput. Second,
the server must deliver the data to the client just-in-time so
that the media data in the client's buffer neither overflows
nor underflows. Third, the latency between the time when the
data is requested and when the data is available for the client
to decode and play back must be short to support interactive
multimedia applications.  To address these challenges, the
conventional wisdom aims to minimize latency and maximize the
utilization of the server's disk.  My thesis research
addressed these challenges
in three novel ways. First, contrary to the conventional wisdom,
I showed that spacing out IOs in time with delay may in fact lead to higher
throughput, since memory, rather than disk, is the resource
bottleneck of the media server \cite{ex49}. I showed how to use memory
judiciously and how to minimize the per-stream cost so that
the throughput can be maximized \cite{ex200}.  Second, I proposed intelligent
data placement and disk scheduling policies to minimize the
initial latency of media data delivery \cite{ex80,ex54}. The analytical model
I developed showed that my schemes can bring the worst-case initial latency
down to just a fraction of a second. 

Finally, I also applied these techniques
to manage parallel disks efficiently \cite{ex90-1}. 
I proposed a resource management scheme that
services requests in {\it cycles} (the time dimension) on
parallel disks (the disk dimension).  Each cycle is divided into {\it
service slots} of equal duration.  An ongoing presentation performs
one IO per cycle, using up one of the slots in that cycle on one of
the disks.  The {\em scheduling policy} assigns IO requests onto
these two dimensional disk-time slots.  It attempts to maintain the
maximum number of free slots open in the near future to service new
requests.  If new requests arrive before the start of a service slot,
they are usually scheduled immediately in the free slots, with little
delay.  If free slots are still available after new requests have
been satisfied, the scheduler assigns existing requests that have the
closest service deadlines to use the remaining slots.  This
effectively defers the deadlines of the most urgent tasks to a later
time, freeing up service slots in the immediate future, as well as
reducing the number of requests that will compete for the free slots
with future new requests.  We coined the
name ``bubble-up'' for this scheme because in it free
slots bubble up to or near the current time.  The simulation results
showed that my scheme can service new requests in under $0.5$ 
second on the average, even when the requests
arrive in a large batch and when the server is near its peak load.
Other schemes can take from 5 seconds to several minutes.        

A {\em data placement} policy is used by the resource management
scheme to balance loads across
disks.  Each presentation is split into {\em chunks} and spread
across the disks.  Like other traditional schemes, my
scheme also replicates popular
presentations.  However, the replication and data distribution are
performed in a way that enhances the run-time bubbling-up of free
slots.  The results showed that two copies
of the popular presentations are sufficient to balance loads and
permit the bubble-up strategy to effectively reduce latencies.

In short, I simulated schemes that manage media servers to
achieve high throughput and low initial latency at low cost.

\subsection{Media Receivers}

As a part of my thesis, I also investigated the receiver-side issues. 
I studied how to manage memory and IOs 
for the receiver \cite{ex100}
to minimize the DRAM requirement for 
cushioning against variability in the
data delivery rate and against differences in delivery and consumption rates.
If one assumes a perfect media server, which can pump out
the data at precisely the rate at which it will be consumed by the receiver,
and one assumes a perfect network, which can deliver the data
at this same rate, then one does not need a buffer at the receiver.
However, the server and network may deliver data at a rate
different from that at which data is consumed, and the rate may not 
be constant, due to network disturbances, traffic congestion, router failures,
server glitches, and so on.  Thus, a receiver cache is necessary 
to handle the delivery-consumption mismatches.
For example, if the delivery rate is smaller than the consumption rate,
the cache must save up enough data before playback starts.
Similarly, the cache can provide playback data to the decoder 
during input lulls.  Furthermore, a receiver cache is also quite 
useful in that it makes the requirements on the server and network
less stringent: the larger the cache, the more flexibility is given 
to the server and network for data delivery.
                                                 
A cache can be made larger by adding disk storage.
A disk cache at the receiver gives the server
much more delivery flexibility.
For instance, a server can then send much more data to the receiver
during a low utilization period (if the network permits),
and less data during a peak utilization period.
In this sense, the receiver assists in server load balancing.
%For example, say, a server with a capacity to service 100
%receivers expects its peak load in 30 minutes.
%If the server is currently servicing $50$ receivers
%that have 60 minutes more of playback time, $50\%$
%of the server's capacity will be tied up during the peak period.
%If the channel capacity is sufficient,
%the server can double the data delivery rate to the current
%$50$ receivers, freeing capacity for the peak period.
%For this to work, the receiver must have a large cache,
%and using a disk for this large cache (as opposed to main memory only)
%significantly reduces costs.
%                               
A large disk cache also makes it possible to deal with
limited channel capacity.
For instance, say a news clip that has a display
rate of $1.5$ Mbps is delivered via a $1.0$ Mbps
channel to the receiver.   Suppose the playback
duration of the news clip is 1000 seconds (about 16 minutes).
During playback, the receiver receives only $1000 \times 1.0$ Mb
but consumes $1000 \times 1.5$ Mb.
Thus, when the clip is started, the receiver needs to have cached
at least $1000 \times 0.5$ Mb, or 62 Mbytes.
Saving this data on disk costs far less than saving it all
in main memory.
               
My contributions in the receiver side study are as follows:
\begin{jlist}
\item
I presented algorithm MEDIC, an adaptive, integrated memory-disk
buffer manager.  The challenge was to develop an algorithm that
\begin{nlist}
\item {\bf 1}. Can efficiently deal with the competing and
unpredictable read and write memory requirements,
\item {\bf 2}. Avoids writing data to disk when it is not necessary,
\item {\bf 3}. Handles out-of-order network data,
\item {\bf 4}. Works on any off-the-shelf disk, and
\item {\bf 5}. Has the right {\it control parameters} for tuning.
\end{nlist}
\item
I derived bounds for the amount of memory a receiver needs
to support given input and playback characteristics.
From these bounds we computed effective settings for
the control parameters, which determine the timing
and sizes of disk IOs.        
\item
I presented a realistic case study, and showed
how input peak rates, disk performance, and other parameters
affect memory and performance.
The results indicate that, thanks to the disk cache,
a relatively modest amount of main memory can support a wide range
of input and media rates.
\end{jlist}

\subsection{Image Databases}

In addition to my thesis work,
I co-designed with my colleagues an image database
for performing content-based
image retrieval. While many papers have been published in this
area, the novelty of our approach lies in two areas. First, we showed that
by using the right features, a similarity search engine
can be tuned to detect unlawful image copies on the world-wide
web \cite{Clw99,Cwl98}. Second, we proposed Clindex \cite{Lch99}, 
an indexing structure that
combines clustering with indexing, to perform approximate similarity
search in high-dimensional spaces. 
We studied how clustering and indexing can be effectively combined.
We experimentally evaluated the Clindex approach,
using a 30,000 image database.
Our results show that Clindex
achieves very high {\em recall}.  That is, it can typically return
more than $90\%$ of what we call the ``golden'' results (i.e., the
best results produced by a linear scan over the entire dataset) with
a few IOs.  We also compared Clindex with other traditional approaches,
to evaluate how effective Clindex's clustering is,
and to understand the gains achievable by
pre-processing data to find clusters.
If data does not have natural clusters, then Clindex will
not improve searches. Fortunately, real datasets rarely occupy a
large-dimensional space uniformly. Our experimental results confirm that,
in the specific case of image data,
significant clusters are formed, and that they can be
exploited to noticeably improve search times.

Theoretical studies~\cite{AM94,IM98,Jk97} have confirmed the cost of
similarity search, independent of the data structure used.  In particular,
it has been shown that if $N$ is the size of a dataset, $D$ is the
dimension, and $D >> log~N$, then no nearest-neighbor algorithm can
be significantly faster than a linear search.
Clindex requires retrieving only a small fraction of the data to
achieve very high recall. We therefore believe that Clindex is a viable
approach for indexing very high-dimensional data for
approximate nearest-neighbor search.


